%%%%%%%%%%%%%%%%%%%%%%%%%%%%%%%%%%%%%%%%%%%%%%%%%%%%%%%%%%
% FILE: chapter01_introduccion.tex
%
% PURPOSE
% -------------------------------------------------------
% Introduces the LaTeX template and explains its
% objectives, philosophy and intended use.
%
%
% ARCHITECTURE ROLE
% -------------------------------------------------------
% First chapter of the template documentation.
%
% Provides conceptual context before explaining
% the technical structure of the project.
%
%
% DEPENDENCIES
% -------------------------------------------------------
% environments.tex
% (definitionbox, notebox, warningbox)
%
%
% NOTES FOR DEVELOPERS
% -------------------------------------------------------
% This chapter should remain conceptual.
%
% Do not include implementation details here.
% Technical explanations belong to later chapters.
%
%%%%%%%%%%%%%%%%%%%%%%%%%%%%%%%%%%%%%%%%%%%%%%%%%%%%%%%%%%


\chapter{Introducción}

Esta plantilla ha sido diseñada para facilitar la redacción de
documentos académicos y técnicos utilizando \LaTeX.

El objetivo principal es proporcionar una estructura modular
que permita separar claramente:

\begin{itemize}
\item configuración del documento
\item estilos tipográficos
\item contenido académico
\item recursos gráficos y bibliográficos
\end{itemize}

Gracias a esta separación, el autor puede concentrarse
exclusivamente en escribir el contenido del documento
sin preocuparse por el formato.



% --------------------------------------------------
% OBJETIVO DE LA PLANTILLA
% --------------------------------------------------

\section{Objetivo de la plantilla}

\begin{definitionbox}[Plantilla académica]
Una plantilla académica es una estructura predefinida
que automatiza la configuración tipográfica y
organizativa de un documento.
\end{definitionbox}

Esta plantilla está diseñada para:

\begin{itemize}
\item trabajos universitarios
\item informes técnicos
\item proyectos de ingeniería
\item Trabajos Fin de Grado (TFG)
\item Trabajos Fin de Máster (TFM)
\end{itemize}



% --------------------------------------------------
% FILOSOFÍA DE DISEÑO
% --------------------------------------------------

\section{Filosofía de diseño}

La plantilla sigue varios principios de diseño inspirados
en prácticas de ingeniería de software.

\begin{itemize}
\item modularidad
\item separación de responsabilidades
\item reutilización
\item documentación interna
\end{itemize}



\begin{notebox}
El objetivo no es únicamente generar un documento
correctamente formateado, sino también demostrar
buenas prácticas en la organización de proyectos
\LaTeX.
\end{notebox}



% --------------------------------------------------
% ORGANIZACIÓN DEL PROYECTO
% --------------------------------------------------

\section{Organización del proyecto}

El proyecto se organiza en varias carpetas que separan
la configuración del documento del contenido académico.



\begin{lstlisting}[caption={Estructura general del proyecto},label={code:estructura}]
main.tex

latex-university-template/
|-- main.tex
|-- 00_config/
|   |-- config.tex
|   `-- packages.tex
|-- 01_frontmatter/
|-- 02_chapters/
|-- 03_figures/
|-- 04_bibliography/
|-- 05_appendices/
|-- 06_styles/
`-- 07_acronyms/
\end{lstlisting}

\source{Plantilla}



% --------------------------------------------------
% USO RECOMENDADO
% --------------------------------------------------

\section{Uso recomendado}

El flujo de trabajo recomendado es el siguiente:

\begin{enumerate}
\item Configurar los metadatos del documento
\item Escribir el contenido en la carpeta \texttt{02\_chapters}
\item Añadir figuras en \texttt{03\_figures}
\item Añadir referencias en \texttt{05\_bibliography}
\end{enumerate}



\begin{warningbox}
Modificar directamente los archivos de configuración
o de estilos sin comprender su función puede producir
errores de compilación o inconsistencias tipográficas.
\end{warningbox}



% --------------------------------------------------
% METODOLOGÍA DE DOCUMENTACIÓN
% --------------------------------------------------

\section{Metodología de documentación}

Todos los archivos de la plantilla siguen una convención
de comentarios estructurados.

Cada archivo comienza con un bloque de documentación
que describe su función dentro del proyecto.

\begin{lstlisting}[caption={Bloque de documentación estándar},label={code:docblock}]
%%%%%%%%%%%%%%%%%%%%%%%%%%%%%%%%%%%%%%%%%%%%%%%%%%%%%%%%%%
% FILE: filename.tex
%
% PURPOSE
% Describe el proposito del archivo
%
% ARCHITECTURE ROLE
% Describe su funcion dentro del sistema
%
% DEPENDENCIES
% Archivos o paquetes necesarios
%
% NOTES FOR DEVELOPERS
% Indicaciones para mantenimiento
%%%%%%%%%%%%%%%%%%%%%%%%%%%%%%%%%%%%%%%%%%%%%%%%%%%%%%%%%%
\end{lstlisting}

\source{Plantilla}



\begin{examplebox}
Este sistema de documentación permite comprender
rápidamente la arquitectura del proyecto al revisar
los archivos de la plantilla.
\end{examplebox}
